\section{Introduction}

Since their invention in the early 80's \cite{fork1981}, femtosecond lasers have rapidly advanced the frontiers of materials science research due to their high peak intensity and outstandingly short temporal resolution.

The ultra-short nature of femtosecond pulses has enabled imaging of ultrafast processes, such as chemical reactions \cite{zewail1994} and the complex energy transfer dynamics between electrons and crystal lattices in solids \cite{shah1996}, whereas the high peak intensity of these pulses allows them to ionise matter and trigger laser-produced plasma (LPP), which emits a variety of spectra, including X-rays \cite{mao2004}.

Studying the dynamics of femtosecond pulses in dielectric media, ranging from the near-UV to the near-infrared, helps clarify the transfer of electromagnetic energy to electrons, followed by electron-lattice interactions that convert energy into heat \cite{mao2004}. However, such studies remain challenging, and some phenomena remain weakly understood for wide-bandgap dielectric materials.

Indeed, these media are transparent under low-intensity irradiation, and when they are irradiated by an intense femtosecond pulse, a cascade of effects begins, ranging from non-linear ionisation occurring within a few femtoseconds to thermal effects occurring on millisecond timescales \cite{bulgakova2010}. Once non-linear ionisation becomes significant, the material loses its transparency, leading to photo-induced damage.

Femtosecond laser processing techniques have an edge over continuous-wave laser machining techniques because of the extremely fast interaction time between the beam and the sample. The damage induced by femtosecond laser pulses is mainly driven by electronic excitation and associated non-linear processes before the material lattice has equilibrated with the excited carriers \cite{mao2004}. This allows for better control over collateral thermal damage and thermal stresses, which are typically associated with continuous laser processing. Also, the non-linear response allows for structural modifications of dielectrics inside the bulk and paves the way for microstructuring that creates subwavelength "voxels" \cite{mao2004}.

The complexity of these nonlinear interactions has spurred extensive research aimed at understanding their fundamental mechanisms and harnessing these properties to develop precision laser processing techniques.

This work presents a simulation of beam propagation in $\mathrm{SiO}_2$ using the nonlinear Schrödinger equation and equations governing electron density evolution. We will explain the physics underlying the model, as well as some of the numerical methods used to solve this complex system.
